%! Departures from Linearity — Unit 2, Lesson 2.5 — AP Practice (Answer Key)
% MIRROR of ../ap_practice/main.tex: each \writelines{n} becomes n terse \ansline{}; the correct MC
% option is wrapped in \ans{}. NO teachernote here — guidance lives in the lesson plan.
% GRADUAL RELEASE: AP-format practice, assigned but NOT scored — the cover's score column reads
% NA for this component. The graded practice is the `homework` component that follows it.
% ONE PAGE (compact budget): two multiple-choice items on a bacteria culture (MC1 is the crux in
% AP clothing — high r, curved residual plot; MC2 is a log-model prediction, DAT-1.J.1), then
% one free-response set on used cars: (a) leverage, (b) influence from output, (c) the SPIRAL —
% interpret the slope in context (Lesson 2.4).
% Page lockstep with ap_practice_key rests on \writelines{n} in the blank matching exactly n
% terse \ansline{} in the key. The next-lesson preview is NOT here: it closes the packet, in
% ../homework/.
\documentclass[12pt]{article}
\usepackage{apstats-article}
\usepackage{apstats-key}   % pulls in -boxes; do NOT also load -boxes

\begin{document}

\pageheader{Unit 2, Lesson 2.5}{AP Practice --- Answer Key}
% NAMESTRIP: no \namedateperiod here — mirrors the blank.

\vspace{0.06in}

% Authored BYTE-IDENTICALLY in ap_practice/ and ap_practice_key/.
\boxguard[8]
\begin{remindbox}
\small
This page is \textbf{not required and not scored} --- your graded homework is the last section
of this packet. Work these AP-format reps for practice and bring questions to study hall.
\end{remindbox}

\vspace{0.06in}

\boxguard[10]
\begin{headlinebox}{goldbg}
\small
\textbf{Section I --- Multiple Choice.} Circle the single best answer.
\end{headlinebox}

\vspace{0.06in}

\begin{enumerate}[leftmargin=1.9em, itemsep=0.10in]

  \item A biologist fits a line to the number of bacteria in a culture (thousands) against
        hours since it was started: $r = 0.96$, but the residual plot shows a clear curve.
        Which conclusion is best?
        \begin{enumerate}[label=\textbf{(\Alph*)}, leftmargin=2.2em, itemsep=1pt]
          \item The line is an appropriate model, because $r$ is close to 1.
          \item The line is an appropriate model, because the residuals add to zero.
          \item \ans{A line is not appropriate: the residuals curve, even though $r$ is high.}
          \item There is no association between hours and the number of bacteria.
          \item The curved residual plot means $r$ was calculated incorrectly.
        \end{enumerate}

  \item The biologist takes the log of each count and fits $\widehat{\log y} = 0.50 + 0.25x$,
        where $x$ is hours and $y$ is bacteria (thousands). The new residual plot shows random
        scatter. What is the predicted number of bacteria (thousands) 4 hours after the start?
        \par\vspace{3pt}\noindent
        \hspace{1.2em}\textbf{(A)} 1.5 \hfill \textbf{(B)} 3.2 \hfill
        \textbf{(C)} 4.5 \hfill \ans{\textbf{(D)} 31.6} \hfill \textbf{(E)} 316

\end{enumerate}

\vspace{0.08in}

\boxguard[10]
\begin{headlinebox}{goldbg}
\small
\textbf{Section II --- Free Response.} Every answer must be written \emph{in the context of the
problem} --- that is what earns credit on the AP exam.
\end{headlinebox}

\vspace{0.06in}

\begin{enumerate}[leftmargin=1.9em, itemsep=0.10in, resume]

  \item \begin{minipage}[t]{0.50\linewidth}\raggedright
        Thirteen used cars of one model: twelve have 22{,}000 to 90{,}000 miles, one has
        210{,}000. Price ($y$, thousands of dollars) against mileage ($x$, thousands of miles):
        \end{minipage}\hfill
        \begin{minipage}[t]{0.46\linewidth}\vspace{0pt}\centering\small
        \setlength{\tabcolsep}{4pt}%
        \renewcommand{\arraystretch}{1.1}
        \begin{tabular}{@{} l l c @{}}
        \rowcolor{sky}
        \textbf{Line fit to} & \textbf{Equation} & $\boldsymbol{r}$ \\
        \hline
        all 13 cars & $\hat{y} = 21.35 - 0.055x$ & $-0.91$ \\
        the 12 others & $\hat{y} = 23.96 - 0.103x$ & $-0.99$ \\
        \end{tabular}
        \end{minipage}

        \vspace{6pt}

        \textbf{(a)} Is the 210{,}000-mile car a high-leverage point? Explain.\\[4pt]
        \ansline{Yes --- 210{,}000 miles is far beyond the other cars' 22{,}000 to 90{,}000.}

        \vspace{4pt}
        \textbf{(b)} Is that car an influential point? Justify with the output.\\[4pt]
        \ansline{Yes --- without it the slope goes from $-0.055$ to $-0.103$ (it nearly doubles)}\\[1pt]
        \ansline{and $r$ goes from $-0.91$ to $-0.99$.}

        \vspace{4pt}
        \textbf{(c)} Interpret the slope of the line for the 12 others in context.\\[4pt]
        \ansline{For each additional 1{,}000 miles, the predicted price of this model drops by}\\[1pt]
        \ansline{about \$103 (0.103 thousand dollars).}

\end{enumerate}

\end{document}
